\documentclass[a4paper,twoside]{article}

\usepackage{morewrites}

\usepackage[svgnames,dvipsnames,x11names]{xcolor}
\usepackage{graphicx}
\usepackage{texsmith-lists}
\usepackage[english]{babel}
\usepackage[colorlinks=true, linkcolor=NavyBlue, urlcolor=NavyBlue, citecolor=NavyBlue]{hyperref}
\usepackage[a4paper]{geometry}
\usepackage{ts-fonts}
\usepackage{booktabs}
\usepackage{tabularx}
\usepackage{float}
\usepackage{seqsplit}
\newcolumntype{Y}{>{\centering\arraybackslash}X}
\newcolumntype{Z}{>{\raggedleft\arraybackslash}X}
\usepackage{ts-code}

\usepackage{lastpage}
\usepackage{refcount}
\makeatletter
\def\ps@plain{
\let\@oddhead\@empty
\let\@evenhead\@empty
\def\@oddfoot{
\hfil
\ifnum\getpagerefnumber{LastPage}>1\relax
\thepage
\fi
\hfil}
\let\@evenfoot\@oddfoot
}
\makeatother

\title{Firmware Review Report\\[0.5em]\large Sensor node, release candidate 1.4.2-rc1}
\author{Firmware Quality Board}\date{March 12, 2026}
\begin{document}
\maketitle
\tableofcontents
\section{Scope}
This report reviews release candidate 1.4.2-rc1 of the sensor node firmware.
Two numbered series run through it: findings raised by the review board, and
the requirements they are checked against. Both are declared in the YAML front
matter under \tscodeinline{press.declare.counters}, so TeXSmith allocates the numbers and both backends ---
\LaTeX{} and Typst --- print the same ones.

The requirement series starts at \tscodeinline{100} because it continues the numbering of
the system specification; the finding series starts at \tscodeinline{1}, the default. The
board considers one issue blocking, finding \hyperref[fw:watchdog]{FW-01}, which is described
further down --- a reference may point forward as well as backward.
\section{Requirements under review}
The four requirements below were selected for this review. The identifier in
the first column is a definition marker, \tscodeinline{\#\{prefix:key\}}: it prints the
formatted number and becomes the anchor every later reference links to.
\begin{center}
\begin{tabularx}{\linewidth}{>{\raggedright\arraybackslash}X>{\raggedright\arraybackslash}X>{\raggedright\arraybackslash}X}
\toprule
\textbf{Id} & \textbf{Requirement} & \textbf{Verification} \\
\midrule
\leavevmode\phantomsection\label{req:watchdog-reset}REQ-100 & The watchdog shall reset the node within two seconds of a stalled main loop. & Fault injection \\
\leavevmode\phantomsection\label{req:ota-rollback}REQ-101 & An interrupted over-the-air update shall leave the previous image bootable. & Power-cut campaign \\
\leavevmode\phantomsection\label{req:key-entropy}REQ-102 & The session key shall be derived from at least 128 bits of hardware entropy. & Code review \\
\leavevmode\phantomsection\label{req:log-retention}REQ-103 & The node shall retain the last 64 log records across a reset. & Manual inspection \\
\bottomrule
\end{tabularx}
\end{center}
Requirement \hyperref[req:key-entropy]{REQ-102} was added after last year's security audit and had
never been verified before this review.
\section{Findings}
\subsection{Summary}
The review raised five findings; the three severe ones are summarised below.
As in the requirement table, the identifiers are defined right in the first
column.
\begin{center}
\begin{tabularx}{\linewidth}{>{\raggedright\arraybackslash}X>{\raggedright\arraybackslash}X>{\raggedright\arraybackslash}X>{\raggedright\arraybackslash}X}
\toprule
\textbf{Id} & \textbf{Severity} & \textbf{Component} & \textbf{Summary} \\
\midrule
\leavevmode\phantomsection\label{fw:watchdog}FW-01 & Blocking & \tscodeinline{hal/watchdog.c} & The watchdog is fed from the I2C completion handler. \\
\leavevmode\phantomsection\label{fw:ota-brick}FW-02 & Blocking & \tscodeinline{ota/apply.c} & A power cut during the swap leaves no bootable image. \\
\leavevmode\phantomsection\label{fw:key-entropy}FW-03 & Major & \tscodeinline{crypto/session.c} & The session key is seeded from the boot counter. \\
\bottomrule
\end{tabularx}
\end{center}
\subsection{Watchdog fed from an interrupt handler}
Feeding the watchdog from the I2C completion handler means the timer keeps
being serviced while the main loop is stalled on a contended bus. Finding
\hyperref[fw:watchdog]{FW-01} therefore violates \hyperref[req:watchdog-reset]{REQ-100} --- on a saturated bus the
node stayed unresponsive for more than forty seconds without ever resetting.
\subsection{Over-the-air update leaves no bootable image}
The updater erases the active slot before validating the incoming image. A
power cut in that window bricks the node, which is exactly what
\hyperref[req:ota-rollback]{REQ-101} forbids; see \hyperref[fw:ota-brick]{FW-02}.
\subsection{Session key seeded from the boot counter}
The seed passed to the key derivation function is the 32-bit boot counter, so a
freshly provisioned node produces a predictable session key. Finding
\hyperref[fw:key-entropy]{FW-03} is raised against \hyperref[req:key-entropy]{REQ-102}.
\subsection{Log buffer wiped on reset}
\label{fw:log-wrap}
The heading above carries a \tscodeinline{\{\#prefix:key\}} attribute, a \emph{silent} definition:
the number appears nowhere in the title, yet the finding is allocated one and
can be referenced like any other. The log ring buffer lives in a \tscodeinline{.bss}
section that the startup code zeroes, so no record survives a reset and
\hyperref[req:log-retention]{REQ-103} is not met. Finding \hyperref[fw:log-wrap]{FW-04} was reported by the field
team before the review started.
\subsection{Additional observation}
Definition markers are ordinary inline constructs and work in running prose
too. While reproducing the watchdog campaign we also noticed \leavevmode\phantomsection\label{fw:rtc-drift}FW-05,
the real-time clock drifting by roughly four seconds a day at 60 °C. No
requirement covers clock accuracy, so this finding is recorded for information
only.
\section{Follow-up}
Findings \hyperref[fw:watchdog]{FW-01} and \hyperref[fw:ota-brick]{FW-02} block the release and must be fixed in
1.4.2-rc2. Finding \hyperref[fw:key-entropy]{FW-03} is scheduled for 1.5.0 together with the
hardware entropy source. The remaining two, \hyperref[fw:log-wrap]{FW-04} and \hyperref[fw:rtc-drift]{FW-05}, are
deferred to the maintenance backlog.

A reference prints the number and nothing else --- \tscodeinline{FW-\allowbreak{}01}, not \tscodeinline{Finding FW-\allowbreak{}01}.
The surrounding noun is written by the author, exactly as for a section
cross-reference; the \tscodeinline{name:} field of the counter is only used in TeXSmith's
diagnostics.
\section{A note on undeclared prefixes}
A marker is substituted only when its prefix is declared in the front matter.
The firmware's Ruby log formatter interpolates values into its message
template: spelled out in prose, it uses \#\{node.id\} for the node identifier and
\#\{sensor:temp\} for the reading. Neither \tscodeinline{node} nor \tscodeinline{sensor} is a declared
counter, so both markers survive into the PDF unchanged --- Ruby and CoffeeScript
interpolations are never mangled, and no warning is emitted. Only \tscodeinline{fw} and
\tscodeinline{req} are declared here, and only those are substituted.
\end{document}
